% Ngân hàng câu hỏi Toán 11 (KNTT) -- Chương 2, Bài 6, Dạng 10
% Dạng: Cấp số cộng: số hạng tổng quát, tổng n số hạng; thực tế
% 100 câu: NB 30 -- TH 40 -- VD 20 -- VDC 10 ; mỗi mức: trắc nghiệm : đúng/sai : trả lời ngắn = 3 : 1 : 1
% Phương án ĐẦU TIÊN trong \cauTN luôn là đáp án đúng (gói trondethi tự đảo).
% Phần II: \y{nội dung}{d|s}. Phần III: đáp án đặt ở đối số cuối.

\cauTN[id=T11-C2-B6-D10-NB-001,muc=NB,loigiai={$u_{25}=u_1+24 d=-67$.}]{Cho cấp số cộng $(u_n)$ có $u_1=5$ và công sai $d=-3$. Số hạng $u_{25}$ bằng}
  {$-67$.}
  {$-68$.}{$-70$.}{$-65$.}

\cauTN[id=T11-C2-B6-D10-NB-002,muc=NB,loigiai={$d=u_2-u_1=-2$.}]{Cho cấp số cộng có các số hạng đầu là $2;\ 0;\ -2;\ -4;\ldots$ Công sai của cấp số cộng là}
  {$-2$.}
  {$-4$.}{$0$.}{$2$.}

\cauTN[id=T11-C2-B6-D10-NB-003,muc=NB,loigiai={Dãy $7;\ 7;\ 7;\ 7$ có hiệu hai số hạng liên tiếp không đổi.}]{Dãy số nào sau đây là một cấp số cộng?}
  {$7;\ 7;\ 7;\ 7$.}
  {$3;\ 5;\ 8;\ 11$.}{$1;\ -1;\ 1;\ -1$.}{$0;\ 1;\ 4;\ 9$.}

\cauTN[id=T11-C2-B6-D10-NB-004,muc=NB,loigiai={$u_{27}=u_1+26 d=61$.}]{Cho cấp số cộng $(u_n)$ có $u_1=9$ và công sai $d=2$. Số hạng $u_{27}$ bằng}
  {$61$.}
  {$63$.}{$60$.}{$54$.}

\cauTN[id=T11-C2-B6-D10-NB-005,muc=NB,loigiai={$d=u_2-u_1=1$.}]{Cho cấp số cộng có các số hạng đầu là $8;\ 9;\ 10;\ 11;\ldots$ Công sai của cấp số cộng là}
  {$1$.}
  {$3$.}{$2$.}{$-1$.}

\cauTN[id=T11-C2-B6-D10-NB-006,muc=NB,loigiai={Dãy $\dfrac12;\ 1;\ \dfrac32;\ 2$ có hiệu hai số hạng liên tiếp không đổi.}]{Dãy số nào sau đây là một cấp số cộng?}
  {$\dfrac12;\ 1;\ \dfrac32;\ 2$.}
  {$1;\ -1;\ 1;\ -1$.}{$1;\ 3;\ 6;\ 10$.}{$3;\ 5;\ 8;\ 11$.}

\cauTN[id=T11-C2-B6-D10-NB-007,muc=NB,loigiai={$u_{15}=u_1+14 d=-94$.}]{Cho cấp số cộng $(u_n)$ có $u_1=-10$ và công sai $d=-6$. Số hạng $u_{15}$ bằng}
  {$-94$.}
  {$-90$.}{$-100$.}{$-95$.}

\cauTN[id=T11-C2-B6-D10-NB-008,muc=NB,loigiai={$d=u_2-u_1=6$.}]{Cho cấp số cộng có các số hạng đầu là $-8;\ -2;\ 4;\ 10;\ldots$ Công sai của cấp số cộng là}
  {$6$.}
  {$8$.}{$7$.}{$5$.}

\cauTN[id=T11-C2-B6-D10-NB-009,muc=NB,loigiai={Dãy $1;\ 4;\ 7;\ 10;\ 13$ có hiệu hai số hạng liên tiếp không đổi.}]{Dãy số nào sau đây là một cấp số cộng?}
  {$1;\ 4;\ 7;\ 10;\ 13$.}
  {$1;\ -1;\ 1;\ -1$.}{$0;\ 1;\ 4;\ 9$.}{$3;\ 5;\ 8;\ 11$.}

\cauTN[id=T11-C2-B6-D10-NB-010,muc=NB,loigiai={$u_{8}=u_1+7 d=-4$.}]{Cho cấp số cộng $(u_n)$ có $u_1=10$ và công sai $d=-2$. Số hạng $u_{8}$ bằng}
  {$-4$.}
  {$-6$.}{$-3$.}{$-16$.}

\cauTN[id=T11-C2-B6-D10-NB-011,muc=NB,loigiai={$d=u_2-u_1=-1$.}]{Cho cấp số cộng có các số hạng đầu là $1;\ 0;\ -1;\ -2;\ldots$ Công sai của cấp số cộng là}
  {$-1$.}
  {$0$.}{$-3$.}{$1$.}

\cauTN[id=T11-C2-B6-D10-NB-012,muc=NB,loigiai={$u_{5}=u_1+4 d=-18$.}]{Cho cấp số cộng $(u_n)$ có $u_1=6$ và công sai $d=-6$. Số hạng $u_{5}$ bằng}
  {$-18$.}
  {$-17$.}{$-16$.}{$-30$.}

\cauTN[id=T11-C2-B6-D10-NB-013,muc=NB,loigiai={$d=u_2-u_1=2$.}]{Cho cấp số cộng có các số hạng đầu là $2;\ 4;\ 6;\ 8;\ldots$ Công sai của cấp số cộng là}
  {$2$.}
  {$-2$.}{$1$.}{$6$.}

\cauTN[id=T11-C2-B6-D10-NB-014,muc=NB,loigiai={Dãy $10;\ 6;\ 2;\ -2$ có hiệu hai số hạng liên tiếp không đổi.}]{Dãy số nào sau đây là một cấp số cộng?}
  {$10;\ 6;\ 2;\ -2$.}
  {$1;\ \dfrac12;\ \dfrac14;\ \dfrac18$.}{$0;\ 1;\ 4;\ 9$.}{$3;\ 5;\ 8;\ 11$.}

\cauTN[id=T11-C2-B6-D10-NB-015,muc=NB,loigiai={$u_{28}=u_1+27 d=-141$.}]{Cho cấp số cộng $(u_n)$ có $u_1=-6$ và công sai $d=-5$. Số hạng $u_{28}$ bằng}
  {$-141$.}
  {$-143$.}{$-139$.}{$-142$.}

\cauTN[id=T11-C2-B6-D10-NB-016,muc=NB,loigiai={$d=u_2-u_1=4$.}]{Cho cấp số cộng có các số hạng đầu là $7;\ 11;\ 15;\ 19;\ldots$ Công sai của cấp số cộng là}
  {$4$.}
  {$2$.}{$6$.}{$5$.}

\cauTN[id=T11-C2-B6-D10-NB-017,muc=NB,loigiai={$u_{14}=u_1+13 d=19$.}]{Cho cấp số cộng $(u_n)$ có $u_1=-7$ và công sai $d=2$. Số hạng $u_{14}$ bằng}
  {$19$.}
  {$21$.}{$17$.}{$28$.}

\cauTN[id=T11-C2-B6-D10-NB-018,muc=NB,loigiai={$d=u_2-u_1=5$.}]{Cho cấp số cộng có các số hạng đầu là $-1;\ 4;\ 9;\ 14;\ldots$ Công sai của cấp số cộng là}
  {$5$.}
  {$4$.}{$-5$.}{$6$.}

\cauDS[id=T11-C2-B6-D10-NB-019,muc=NB,loigiai={$u_n=-8-4(n-1)$; $u_{7}=-32$; $S_{12}=-360$; $d<0$. a) Đúng. b) Đúng. c) Đúng. d) Đúng.}]{Cho cấp số cộng $(u_n)$ có $u_1=-8$ và $d=-4$.}{%
  \y{$u_2=-12$.}{d}
  \y{$u_{7}=-32$.}{d}
  \y{Tổng $12$ số hạng đầu bằng $-360$.}{d}
  \y{Cấp số cộng là dãy số giảm.}{d}}

\cauDS[id=T11-C2-B6-D10-NB-020,muc=NB,loigiai={$u_n=1-2(n-1)$; $u_{9}=-15$; $S_{14}=-168$; $d<0$. a) Đúng. b) Sai. c) Đúng. d) Đúng.}]{Cho cấp số cộng $(u_n)$ có $u_1=1$ và $d=-2$.}{%
  \y{$u_2=-1$.}{d}
  \y{$u_{9}=-17$.}{s}
  \y{Tổng $14$ số hạng đầu bằng $-168$.}{d}
  \y{Cấp số cộng là dãy số giảm.}{d}}

\cauDS[id=T11-C2-B6-D10-NB-021,muc=NB,loigiai={$u_n=-1+5(n-1)$; $u_{14}=64$; $S_{5}=45$; $d>0$. a) Đúng. b) Sai. c) Sai. d) Sai.}]{Cho cấp số cộng $(u_n)$ có $u_1=-1$ và $d=5$.}{%
  \y{$u_2=4$.}{d}
  \y{$u_{14}=69$.}{s}
  \y{Tổng $5$ số hạng đầu bằng $50$.}{s}
  \y{Cấp số cộng là dãy số giảm.}{s}}

\cauDS[id=T11-C2-B6-D10-NB-022,muc=NB,loigiai={$u_n=-5-3(n-1)$; $u_{10}=-32$; $S_{12}=-258$; $d<0$. a) Đúng. b) Sai. c) Sai. d) Đúng.}]{Cho cấp số cộng $(u_n)$ có $u_1=-5$ và $d=-3$.}{%
  \y{$u_2=-8$.}{d}
  \y{$u_{10}=-35$.}{s}
  \y{Tổng $12$ số hạng đầu bằng $-261$.}{s}
  \y{Cấp số cộng là dãy số giảm.}{d}}

\cauDS[id=T11-C2-B6-D10-NB-023,muc=NB,loigiai={$u_n=5+5(n-1)$; $u_{13}=65$; $S_{5}=75$; $d>0$. a) Đúng. b) Đúng. c) Đúng. d) Sai.}]{Cho cấp số cộng $(u_n)$ có $u_1=5$ và $d=5$.}{%
  \y{$u_2=10$.}{d}
  \y{$u_{13}=65$.}{d}
  \y{Tổng $5$ số hạng đầu bằng $75$.}{d}
  \y{Cấp số cộng là dãy số giảm.}{s}}

\cauDS[id=T11-C2-B6-D10-NB-024,muc=NB,loigiai={$u_n=9-4(n-1)$; $u_{11}=-31$; $S_{9}=-63$; $d<0$. a) Đúng. b) Sai. c) Đúng. d) Sai.}]{Cho cấp số cộng $(u_n)$ có $u_1=9$ và $d=-4$.}{%
  \y{$u_2=5$.}{d}
  \y{$u_{11}=-35$.}{s}
  \y{Tổng $9$ số hạng đầu bằng $-63$.}{d}
  \y{Cấp số cộng là dãy số tăng.}{s}}

\cauTLN[id=T11-C2-B6-D10-NB-025,muc=NB,loigiai={$S_{22}=\dfrac{22[2u_1+21 d]}{2}=946$.}]{Cho cấp số cộng $(u_n)$ có $u_1=1$, $d=4$. Tính tổng $22$ số hạng đầu tiên.}{946}

\cauTLN[id=T11-C2-B6-D10-NB-026,muc=NB,loigiai={$S_{9}=\dfrac{9[2u_1+8 d]}{2}=180$.}]{Cho cấp số cộng $(u_n)$ có $u_1=8$, $d=3$. Tính tổng $9$ số hạng đầu tiên.}{180}

\cauTLN[id=T11-C2-B6-D10-NB-027,muc=NB,loigiai={$S_{24}=\dfrac{24[2u_1+23 d]}{2}=348$.}]{Cho cấp số cộng $(u_n)$ có $u_1=3$, $d=1$. Tính tổng $24$ số hạng đầu tiên.}{348}

\cauTLN[id=T11-C2-B6-D10-NB-028,muc=NB,loigiai={$S_{21}=\dfrac{21[2u_1+20 d]}{2}=63$.}]{Cho cấp số cộng $(u_n)$ có $u_1=-7$, $d=1$. Tính tổng $21$ số hạng đầu tiên.}{63}

\cauTLN[id=T11-C2-B6-D10-NB-029,muc=NB,loigiai={$S_{19}=\dfrac{19[2u_1+18 d]}{2}=703$.}]{Cho cấp số cộng $(u_n)$ có $u_1=10$, $d=3$. Tính tổng $19$ số hạng đầu tiên.}{703}

\cauTLN[id=T11-C2-B6-D10-NB-030,muc=NB,loigiai={$S_{20}=\dfrac{20[2u_1+19 d]}{2}=-150$.}]{Cho cấp số cộng $(u_n)$ có $u_1=2$, $d=-1$. Tính tổng $20$ số hạng đầu tiên.}{-150}

\cauTN[id=T11-C2-B6-D10-TH-031,muc=TH,loigiai={Giải hệ được $u_1=7$, $d=-2$; $u_{29}=-49$.}]{Cho cấp số cộng $(u_n)$ có $u_{7}=-5$ và $u_{13}=-17$. Số hạng $u_{29}$ bằng}
  {$-49$.}
  {$-50$.}{$-51$.}{$-48$.}

\cauTN[id=T11-C2-B6-D10-TH-032,muc=TH,loigiai={$u_1=9$, $d=-3$, $S_{5}=15$.}]{Cho cấp số cộng $(u_n)$ có $u_{4}=0$ và $u_{8}=-12$. Tổng $5$ số hạng đầu tiên của cấp số cộng bằng}
  {$15$.}
  {$17$.}{$-1$.}{$-15$.}

\cauTN[id=T11-C2-B6-D10-TH-033,muc=TH,loigiai={$u_{23}=u_1+22 d=47$.}]{Cho cấp số cộng $(u_n)$ có $u_1=3$ và công sai $d=2$. Số hạng $u_{23}$ bằng}
  {$47$.}
  {$48$.}{$46$.}{$49$.}

\cauTN[id=T11-C2-B6-D10-TH-034,muc=TH,loigiai={Giải hệ được $u_1=0$, $d=-3$; $u_{29}=-84$.}]{Cho cấp số cộng $(u_n)$ có $u_{2}=-3$ và $u_{5}=-12$. Số hạng $u_{29}$ bằng}
  {$-84$.}
  {$-82$.}{$-85$.}{$-81$.}

\cauTN[id=T11-C2-B6-D10-TH-035,muc=TH,loigiai={$u_1=-1$, $d=-3$, $S_{7}=-70$.}]{Cho cấp số cộng $(u_n)$ có $u_{8}=-22$ và $u_{2}=-4$. Tổng $7$ số hạng đầu tiên của cấp số cộng bằng}
  {$-70$.}
  {$-71$.}{$-68$.}{$-92$.}

\cauTN[id=T11-C2-B6-D10-TH-036,muc=TH,loigiai={$u_{18}=u_1+17 d=-59$.}]{Cho cấp số cộng $(u_n)$ có $u_1=-8$ và công sai $d=-3$. Số hạng $u_{18}$ bằng}
  {$-59$.}
  {$-54$.}{$-56$.}{$-62$.}

\cauTN[id=T11-C2-B6-D10-TH-037,muc=TH,loigiai={Giải hệ được $u_1=-10$, $d=1$; $u_{28}=17$.}]{Cho cấp số cộng $(u_n)$ có $u_{11}=0$ và $u_{14}=3$. Số hạng $u_{28}$ bằng}
  {$17$.}
  {$15$.}{$16$.}{$18$.}

\cauTN[id=T11-C2-B6-D10-TH-038,muc=TH,loigiai={$u_1=-10$, $d=4$, $S_{19}=494$.}]{Cho cấp số cộng $(u_n)$ có $u_{9}=22$ và $u_{3}=-2$. Tổng $19$ số hạng đầu tiên của cấp số cộng bằng}
  {$494$.}
  {$1178$.}{$495$.}{$493$.}

\cauTN[id=T11-C2-B6-D10-TH-039,muc=TH,loigiai={$u_{24}=u_1+23 d=129$.}]{Cho cấp số cộng $(u_n)$ có $u_1=-9$ và công sai $d=6$. Số hạng $u_{24}$ bằng}
  {$129$.}
  {$131$.}{$135$.}{$128$.}

\cauTN[id=T11-C2-B6-D10-TH-040,muc=TH,loigiai={Giải hệ được $u_1=-8$, $d=-2$; $u_{23}=-52$.}]{Cho cấp số cộng $(u_n)$ có $u_{11}=-28$ và $u_{10}=-26$. Số hạng $u_{23}$ bằng}
  {$-52$.}
  {$-53$.}{$-54$.}{$-51$.}

\cauTN[id=T11-C2-B6-D10-TH-041,muc=TH,loigiai={$u_1=-6$, $d=-2$, $S_{11}=-176$.}]{Cho cấp số cộng $(u_n)$ có $u_{4}=-12$ và $u_{6}=-16$. Tổng $11$ số hạng đầu tiên của cấp số cộng bằng}
  {$-176$.}
  {$-175$.}{$-174$.}{$-286$.}

\cauTN[id=T11-C2-B6-D10-TH-042,muc=TH,loigiai={$u_{27}=u_1+26 d=125$.}]{Cho cấp số cộng $(u_n)$ có $u_1=-5$ và công sai $d=5$. Số hạng $u_{27}$ bằng}
  {$125$.}
  {$126$.}{$127$.}{$123$.}

\cauTN[id=T11-C2-B6-D10-TH-043,muc=TH,loigiai={Giải hệ được $u_1=-7$, $d=5$; $u_{16}=68$.}]{Cho cấp số cộng $(u_n)$ có $u_{10}=38$ và $u_{3}=3$. Số hạng $u_{16}$ bằng}
  {$68$.}
  {$70$.}{$73$.}{$67$.}

\cauTN[id=T11-C2-B6-D10-TH-044,muc=TH,loigiai={$u_1=7$, $d=1$, $S_{15}=210$.}]{Cho cấp số cộng $(u_n)$ có $u_{8}=14$ và $u_{9}=15$. Tổng $15$ số hạng đầu tiên của cấp số cộng bằng}
  {$210$.}
  {$315$.}{$209$.}{$224$.}

\cauTN[id=T11-C2-B6-D10-TH-045,muc=TH,loigiai={$u_{23}=u_1+22 d=-106$.}]{Cho cấp số cộng $(u_n)$ có $u_1=4$ và công sai $d=-5$. Số hạng $u_{23}$ bằng}
  {$-106$.}
  {$-111$.}{$-107$.}{$-115$.}

\cauTN[id=T11-C2-B6-D10-TH-046,muc=TH,loigiai={Giải hệ được $u_1=2$, $d=-5$; $u_{22}=-103$.}]{Cho cấp số cộng $(u_n)$ có $u_{12}=-53$ và $u_{10}=-43$. Số hạng $u_{22}$ bằng}
  {$-103$.}
  {$-101$.}{$-105$.}{$-104$.}

\cauTN[id=T11-C2-B6-D10-TH-047,muc=TH,loigiai={$u_{27}=u_1+26 d=-52$.}]{Cho cấp số cộng $(u_n)$ có $u_1=0$ và công sai $d=-2$. Số hạng $u_{27}$ bằng}
  {$-52$.}
  {$-54$.}{$-51$.}{$-53$.}

\cauTN[id=T11-C2-B6-D10-TH-048,muc=TH,loigiai={Giải hệ được $u_1=-10$, $d=5$; $u_{19}=80$.}]{Cho cấp số cộng $(u_n)$ có $u_{7}=20$ và $u_{4}=5$. Số hạng $u_{19}$ bằng}
  {$80$.}
  {$82$.}{$79$.}{$81$.}

\cauTN[id=T11-C2-B6-D10-TH-049,muc=TH,loigiai={$u_1=-4$, $d=4$, $S_{13}=260$.}]{Cho cấp số cộng $(u_n)$ có $u_{4}=8$ và $u_{9}=28$. Tổng $13$ số hạng đầu tiên của cấp số cộng bằng}
  {$260$.}
  {$312$.}{$264$.}{$261$.}

\cauTN[id=T11-C2-B6-D10-TH-050,muc=TH,loigiai={$u_{21}=u_1+20 d=66$.}]{Cho cấp số cộng $(u_n)$ có $u_1=6$ và công sai $d=3$. Số hạng $u_{21}$ bằng}
  {$66$.}
  {$68$.}{$67$.}{$64$.}

\cauTN[id=T11-C2-B6-D10-TH-051,muc=TH,loigiai={Giải hệ được $u_1=-6$, $d=6$; $u_{29}=162$.}]{Cho cấp số cộng $(u_n)$ có $u_{2}=0$ và $u_{4}=12$. Số hạng $u_{29}$ bằng}
  {$162$.}
  {$168$.}{$161$.}{$156$.}

\cauTN[id=T11-C2-B6-D10-TH-052,muc=TH,loigiai={$u_{7}=u_1+6 d=-32$.}]{Cho cấp số cộng $(u_n)$ có $u_1=4$ và công sai $d=-6$. Số hạng $u_{7}$ bằng}
  {$-32$.}
  {$-31$.}{$-30$.}{$-26$.}

\cauTN[id=T11-C2-B6-D10-TH-053,muc=TH,loigiai={Giải hệ được $u_1=10$, $d=5$; $u_{17}=90$.}]{Cho cấp số cộng $(u_n)$ có $u_{14}=75$ và $u_{11}=60$. Số hạng $u_{17}$ bằng}
  {$90$.}
  {$89$.}{$85$.}{$91$.}

\cauTN[id=T11-C2-B6-D10-TH-054,muc=TH,loigiai={$u_1=-8$, $d=4$, $S_{7}=28$.}]{Cho cấp số cộng $(u_n)$ có $u_{4}=4$ và $u_{2}=-4$. Tổng $7$ số hạng đầu tiên của cấp số cộng bằng}
  {$28$.}
  {$36$.}{$26$.}{$27$.}

\cauDS[id=T11-C2-B6-D10-TH-055,muc=TH,loigiai={$u_n=-1+5(n-1)$; $u_{9}=39$; $S_{5}=45$; $d>0$. a) Đúng. b) Đúng. c) Đúng. d) Sai.}]{Cho cấp số cộng $(u_n)$ có $u_1=-1$ và $d=5$.}{%
  \y{$u_2=4$.}{d}
  \y{$u_{9}=39$.}{d}
  \y{Tổng $5$ số hạng đầu bằng $45$.}{d}
  \y{Cấp số cộng là dãy số giảm.}{s}}

\cauDS[id=T11-C2-B6-D10-TH-056,muc=TH,loigiai={$u_n=-5+5(n-1)$; $u_{11}=45$; $S_{6}=45$; $d>0$. a) Đúng. b) Đúng. c) Sai. d) Đúng.}]{Cho cấp số cộng $(u_n)$ có $u_1=-5$ và $d=5$.}{%
  \y{$u_2=0$.}{d}
  \y{$u_{11}=45$.}{d}
  \y{Tổng $6$ số hạng đầu bằng $50$.}{s}
  \y{Cấp số cộng là dãy số tăng.}{d}}

\cauDS[id=T11-C2-B6-D10-TH-057,muc=TH,loigiai={$u_n=-1+(n-1)$; $u_{10}=8$; $S_{14}=77$; $d>0$. a) Đúng. b) Đúng. c) Đúng. d) Sai.}]{Cho cấp số cộng $(u_n)$ có $u_1=-1$ và $d=1$.}{%
  \y{$u_2=0$.}{d}
  \y{$u_{10}=8$.}{d}
  \y{Tổng $14$ số hạng đầu bằng $77$.}{d}
  \y{Cấp số cộng là dãy số giảm.}{s}}

\cauDS[id=T11-C2-B6-D10-TH-058,muc=TH,loigiai={$u_n=-6+(n-1)$; $u_{9}=2$; $S_{9}=-18$; $d>0$. a) Đúng. b) Sai. c) Đúng. d) Đúng.}]{Cho cấp số cộng $(u_n)$ có $u_1=-6$ và $d=1$.}{%
  \y{$u_2=-5$.}{d}
  \y{$u_{9}=3$.}{s}
  \y{Tổng $9$ số hạng đầu bằng $-18$.}{d}
  \y{Cấp số cộng là dãy số tăng.}{d}}

\cauDS[id=T11-C2-B6-D10-TH-059,muc=TH,loigiai={$u_n=9-5(n-1)$; $u_{14}=-56$; $S_{12}=-222$; $d<0$. a) Đúng. b) Sai. c) Đúng. d) Sai.}]{Cho cấp số cộng $(u_n)$ có $u_1=9$ và $d=-5$.}{%
  \y{$u_2=4$.}{d}
  \y{$u_{14}=-61$.}{s}
  \y{Tổng $12$ số hạng đầu bằng $-222$.}{d}
  \y{Cấp số cộng là dãy số tăng.}{s}}

\cauDS[id=T11-C2-B6-D10-TH-060,muc=TH,loigiai={$u_n=-1-4(n-1)$; $u_{7}=-25$; $S_{5}=-45$; $d<0$. a) Đúng. b) Sai. c) Đúng. d) Sai.}]{Cho cấp số cộng $(u_n)$ có $u_1=-1$ và $d=-4$.}{%
  \y{$u_2=-5$.}{d}
  \y{$u_{7}=-29$.}{s}
  \y{Tổng $5$ số hạng đầu bằng $-45$.}{d}
  \y{Cấp số cộng là dãy số tăng.}{s}}

\cauDS[id=T11-C2-B6-D10-TH-061,muc=TH,loigiai={$u_n=-4+5(n-1)$; $u_{12}=51$; $S_{6}=51$; $d>0$. a) Đúng. b) Sai. c) Sai. d) Sai.}]{Cho cấp số cộng $(u_n)$ có $u_1=-4$ và $d=5$.}{%
  \y{$u_2=1$.}{d}
  \y{$u_{12}=56$.}{s}
  \y{Tổng $6$ số hạng đầu bằng $56$.}{s}
  \y{Cấp số cộng là dãy số giảm.}{s}}

\cauDS[id=T11-C2-B6-D10-TH-062,muc=TH,loigiai={$u_n=3+3(n-1)$; $u_{11}=33$; $S_{11}=198$; $d>0$. a) Đúng. b) Sai. c) Đúng. d) Sai.}]{Cho cấp số cộng $(u_n)$ có $u_1=3$ và $d=3$.}{%
  \y{$u_2=6$.}{d}
  \y{$u_{11}=36$.}{s}
  \y{Tổng $11$ số hạng đầu bằng $198$.}{d}
  \y{Cấp số cộng là dãy số giảm.}{s}}

\cauTLN[id=T11-C2-B6-D10-TH-063,muc=TH,loigiai={Giải $\dfrac{n[2\cdot(6)+(n-1)\cdot3]}{2}=756$, $n\in\mathbb{N}^*$ được $n=21$.}]{Cho cấp số cộng $(u_n)$ có $u_1=6$, $d=3$. Biết tổng $n$ số hạng đầu bằng $756$. Tìm $n$.}{21}

\cauTLN[id=T11-C2-B6-D10-TH-064,muc=TH,loigiai={$S_{6}=\dfrac{6[2u_1+5 d]}{2}=-45$.}]{Cho cấp số cộng $(u_n)$ có $u_1=-10$, $d=1$. Tính tổng $6$ số hạng đầu tiên.}{-45}

\cauTLN[id=T11-C2-B6-D10-TH-065,muc=TH,loigiai={Giải $\dfrac{n[2\cdot(3)+(n-1)\cdot4]}{2}=528$, $n\in\mathbb{N}^*$ được $n=16$.}]{Cho cấp số cộng $(u_n)$ có $u_1=3$, $d=4$. Biết tổng $n$ số hạng đầu bằng $528$. Tìm $n$.}{16}

\cauTLN[id=T11-C2-B6-D10-TH-066,muc=TH,loigiai={Giải $\dfrac{n[2\cdot(-4)+(n-1)\cdot4]}{2}=756$, $n\in\mathbb{N}^*$ được $n=21$.}]{Cho cấp số cộng $(u_n)$ có $u_1=-4$, $d=4$. Biết tổng $n$ số hạng đầu bằng $756$. Tìm $n$.}{21}

\cauTLN[id=T11-C2-B6-D10-TH-067,muc=TH,loigiai={$S_{16}=\dfrac{16[2u_1+15 d]}{2}=-632$.}]{Cho cấp số cộng $(u_n)$ có $u_1=-2$, $d=-5$. Tính tổng $16$ số hạng đầu tiên.}{-632}

\cauTLN[id=T11-C2-B6-D10-TH-068,muc=TH,loigiai={Giải $\dfrac{n[2\cdot(3)+(n-1)\cdot6]}{2}=1728$, $n\in\mathbb{N}^*$ được $n=24$.}]{Cho cấp số cộng $(u_n)$ có $u_1=3$, $d=6$. Biết tổng $n$ số hạng đầu bằng $1728$. Tìm $n$.}{24}

\cauTLN[id=T11-C2-B6-D10-TH-069,muc=TH,loigiai={$S_{11}=\dfrac{11[2u_1+10 d]}{2}=-44$.}]{Cho cấp số cộng $(u_n)$ có $u_1=1$, $d=-1$. Tính tổng $11$ số hạng đầu tiên.}{-44}

\cauTLN[id=T11-C2-B6-D10-TH-070,muc=TH,loigiai={Giải $\dfrac{n[2\cdot(-3)+(n-1)\cdot2]}{2}=252$, $n\in\mathbb{N}^*$ được $n=18$.}]{Cho cấp số cộng $(u_n)$ có $u_1=-3$, $d=2$. Biết tổng $n$ số hạng đầu bằng $252$. Tìm $n$.}{18}

\cauTN[id=T11-C2-B6-D10-VD-071,muc=VD,loigiai={$S_{18}=\dfrac{18(2\cdot17+17\cdot1)}{2}=459$ ghế.}]{Khu vực ghế khán giả của rạp chiếu phim của một trung tâm thương mại được bố trí theo hình rẻ quạt gồm $18$ hàng ghế. Hàng ghế đầu tiên (gần sân khấu nhất) có $17$ ghế, mỗi hàng ghế phía sau nhiều hơn hàng liền trước đúng $1$ ghế. Tổng số ghế của khu vực khán giả là}
  {$459$.}
  {$460$.}{$630$.}{$458$.}

\cauTN[id=T11-C2-B6-D10-VD-072,muc=VD,loigiai={Lương tháng các năm lập thành cấp số cộng $u_1=7$, $d=\dfrac{1}{2}$; tổng $=12S_{5}=480$ triệu đồng.}]{Một doanh nghiệp may mặc ở Bình Dương tuyển kĩ sư với mức lương năm đầu tiên là $7$ triệu đồng/tháng. Theo hợp đồng, sau mỗi năm làm việc, mức lương mỗi tháng được tăng thêm $\dfrac{1}{2}$ triệu đồng so với năm liền trước. Tổng số tiền lương (triệu đồng) kĩ sư nhận được sau $5$ năm làm việc là}
  {$480$.}
  {$482$.}{$479$.}{$570$.}

\cauTN[id=T11-C2-B6-D10-VD-073,muc=VD,loigiai={Giải hệ được $u_1=-7$, $d=-2$; $u_{19}=-43$.}]{Cho cấp số cộng $(u_n)$ có $u_{4}=-13$ và $u_{10}=-25$. Số hạng $u_{19}$ bằng}
  {$-43$.}
  {$-45$.}{$-42$.}{$-41$.}

\cauTN[id=T11-C2-B6-D10-VD-074,muc=VD,loigiai={Số ghế các hàng lập thành cấp số cộng với $u_1=24$, $d=1$; $u_{18}=41$.}]{Khu vực ghế khán giả của nhà văn hoá thiếu nhi của tỉnh được bố trí theo hình rẻ quạt gồm $18$ hàng ghế. Hàng ghế đầu tiên (gần sân khấu nhất) có $24$ ghế, mỗi hàng ghế phía sau nhiều hơn hàng liền trước đúng $1$ ghế. Số ghế ở hàng cuối cùng là}
  {$41$.}
  {$40$.}{$39$.}{$42$.}

\cauTN[id=T11-C2-B6-D10-VD-075,muc=VD,loigiai={Giải hệ được $u_1=-4$, $d=-1$; $u_{26}=-29$.}]{Cho cấp số cộng $(u_n)$ có $u_{2}=-5$ và $u_{3}=-6$. Số hạng $u_{26}$ bằng}
  {$-29$.}
  {$-31$.}{$-28$.}{$-30$.}

\cauTN[id=T11-C2-B6-D10-VD-076,muc=VD,loigiai={$S_{21}=\dfrac{21(2\cdot12+20\cdot1)}{2}=462$ ghế.}]{Khu vực ghế khán giả của rạp chiếu phim của một trung tâm thương mại được bố trí theo hình rẻ quạt gồm $21$ hàng ghế. Hàng ghế đầu tiên (gần sân khấu nhất) có $12$ ghế, mỗi hàng ghế phía sau nhiều hơn hàng liền trước đúng $1$ ghế. Tổng số ghế của khu vực khán giả là}
  {$462$.}
  {$460$.}{$693$.}{$463$.}

\cauTN[id=T11-C2-B6-D10-VD-077,muc=VD,loigiai={Lương tháng các năm lập thành cấp số cộng $u_1=8$, $d=\dfrac{2}{5}$; tổng $=12S_{5}=528$ triệu đồng.}]{Công ty TNHH Hoàng Long tuyển kĩ sư với mức lương năm đầu tiên là $8$ triệu đồng/tháng. Theo hợp đồng, sau mỗi năm làm việc, mức lương mỗi tháng được tăng thêm $\dfrac{2}{5}$ triệu đồng so với năm liền trước. Tổng số tiền lương (triệu đồng) kĩ sư nhận được sau $5$ năm làm việc là}
  {$528$.}
  {$529$.}{$600$.}{$526$.}

\cauTN[id=T11-C2-B6-D10-VD-078,muc=VD,loigiai={Giải hệ được $u_1=-4$, $d=2$; $u_{23}=40$.}]{Cho cấp số cộng $(u_n)$ có $u_{4}=2$ và $u_{10}=14$. Số hạng $u_{23}$ bằng}
  {$40$.}
  {$38$.}{$42$.}{$39$.}

\cauTN[id=T11-C2-B6-D10-VD-079,muc=VD,loigiai={Số ghế các hàng lập thành cấp số cộng với $u_1=20$, $d=4$; $u_{14}=72$.}]{Khu vực ghế khán giả của Nhà hát Lớn Hà Nội được bố trí theo hình rẻ quạt gồm $14$ hàng ghế. Hàng ghế đầu tiên (gần sân khấu nhất) có $20$ ghế, mỗi hàng ghế phía sau nhiều hơn hàng liền trước đúng $4$ ghế. Số ghế ở hàng cuối cùng là}
  {$72$.}
  {$68$.}{$71$.}{$70$.}

\cauTN[id=T11-C2-B6-D10-VD-080,muc=VD,loigiai={Giải hệ được $u_1=-5$, $d=-5$; $u_{29}=-145$.}]{Cho cấp số cộng $(u_n)$ có $u_{11}=-55$ và $u_{2}=-10$. Số hạng $u_{29}$ bằng}
  {$-145$.}
  {$-143$.}{$-144$.}{$-150$.}

\cauTN[id=T11-C2-B6-D10-VD-081,muc=VD,loigiai={Số ghế các hàng lập thành cấp số cộng với $u_1=23$, $d=1$; $u_{18}=40$.}]{Khu vực ghế khán giả của rạp chiếu phim của một trung tâm thương mại được bố trí theo hình rẻ quạt gồm $18$ hàng ghế. Hàng ghế đầu tiên (gần sân khấu nhất) có $23$ ghế, mỗi hàng ghế phía sau nhiều hơn hàng liền trước đúng $1$ ghế. Số ghế ở hàng cuối cùng là}
  {$40$.}
  {$42$.}{$39$.}{$41$.}

\cauTN[id=T11-C2-B6-D10-VD-082,muc=VD,loigiai={Giải hệ được $u_1=-5$, $d=6$; $u_{27}=151$.}]{Cho cấp số cộng $(u_n)$ có $u_{5}=19$ và $u_{6}=25$. Số hạng $u_{27}$ bằng}
  {$151$.}
  {$152$.}{$149$.}{$157$.}

\cauDS[id=T11-C2-B6-D10-VD-083,muc=VD,loigiai={Cấp số cộng $u_1=12$, $d=2$: $u_{3}=16$, $u_{11}=32$, $S_{11}=242$. a) Đúng. b) Sai. c) Sai. d) Đúng.}]{Khu vực ghế khán giả của nhà văn hoá thiếu nhi của tỉnh được bố trí theo hình rẻ quạt gồm $11$ hàng ghế. Hàng ghế đầu tiên (gần sân khấu nhất) có $12$ ghế, mỗi hàng ghế phía sau nhiều hơn hàng liền trước đúng $2$ ghế.}{%
  \y{Hàng thứ hai có $14$ ghế.}{d}
  \y{Hàng thứ $3$ có $18$ ghế.}{s}
  \y{Hàng cuối cùng có $34$ ghế.}{s}
  \y{Khu vực khán giả có tất cả $242$ ghế.}{d}}

\cauDS[id=T11-C2-B6-D10-VD-084,muc=VD,loigiai={Cấp số cộng $u_1=23$, $d=3$: $u_{9}=47$, $u_{26}=98$, $S_{26}=1573$. a) Đúng. b) Đúng. c) Đúng. d) Đúng.}]{Khu vực ghế khán giả của rạp chiếu phim của một trung tâm thương mại được bố trí theo hình rẻ quạt gồm $26$ hàng ghế. Hàng ghế đầu tiên (gần sân khấu nhất) có $23$ ghế, mỗi hàng ghế phía sau nhiều hơn hàng liền trước đúng $3$ ghế.}{%
  \y{Hàng thứ hai có $26$ ghế.}{d}
  \y{Hàng thứ $9$ có $47$ ghế.}{d}
  \y{Hàng cuối cùng có $98$ ghế.}{d}
  \y{Khu vực khán giả có tất cả $1573$ ghế.}{d}}

\cauDS[id=T11-C2-B6-D10-VD-085,muc=VD,loigiai={Cấp số cộng $u_1=13$, $d=2$: $u_{4}=19$, $u_{17}=45$, $S_{17}=493$. a) Đúng. b) Đúng. c) Đúng. d) Sai.}]{Khu vực ghế khán giả của rạp chiếu phim của một trung tâm thương mại được bố trí theo hình rẻ quạt gồm $17$ hàng ghế. Hàng ghế đầu tiên (gần sân khấu nhất) có $13$ ghế, mỗi hàng ghế phía sau nhiều hơn hàng liền trước đúng $2$ ghế.}{%
  \y{Hàng thứ hai có $15$ ghế.}{d}
  \y{Hàng thứ $4$ có $19$ ghế.}{d}
  \y{Hàng cuối cùng có $45$ ghế.}{d}
  \y{Khu vực khán giả có tất cả $506$ ghế.}{s}}

\cauDS[id=T11-C2-B6-D10-VD-086,muc=VD,loigiai={Cấp số cộng $u_1=21$, $d=1$: $u_{7}=27$, $u_{12}=32$, $S_{12}=318$. a) Đúng. b) Sai. c) Sai. d) Đúng.}]{Khu vực ghế khán giả của rạp chiếu phim của một trung tâm thương mại được bố trí theo hình rẻ quạt gồm $12$ hàng ghế. Hàng ghế đầu tiên (gần sân khấu nhất) có $21$ ghế, mỗi hàng ghế phía sau nhiều hơn hàng liền trước đúng $1$ ghế.}{%
  \y{Hàng thứ hai có $22$ ghế.}{d}
  \y{Hàng thứ $7$ có $28$ ghế.}{s}
  \y{Hàng cuối cùng có $33$ ghế.}{s}
  \y{Khu vực khán giả có tất cả $318$ ghế.}{d}}

\cauTLN[id=T11-C2-B6-D10-VD-087,muc=VD,loigiai={$S_{26}=1690$.}]{Khu vực ghế khán giả của rạp chiếu phim của một trung tâm thương mại được bố trí theo hình rẻ quạt gồm $26$ hàng ghế. Hàng ghế đầu tiên (gần sân khấu nhất) có $15$ ghế, mỗi hàng ghế phía sau nhiều hơn hàng liền trước đúng $4$ ghế. Tính tổng số ghế của khu vực khán giả.}{1690}

\cauTLN[id=T11-C2-B6-D10-VD-088,muc=VD,loigiai={Tìm $n$ lớn nhất với $\dfrac{n[2\cdot3+(n-1)\cdot3]}{2}\le191$, được $n=10$.}]{Người ta xếp các thùng sơn thành một khối hình thang trong kho: hàng trên cùng có $3$ thùng, mỗi hàng bên dưới nhiều hơn hàng liền trên $3$ thùng. Có $191$ thùng sơn. Hỏi xếp được nhiều nhất bao nhiêu hàng đầy đủ theo quy luật trên?}{10}

\cauTLN[id=T11-C2-B6-D10-VD-089,muc=VD,loigiai={$S_{26}=1261$.}]{Khu vực ghế khán giả của Nhà hát Lớn Hà Nội được bố trí theo hình rẻ quạt gồm $26$ hàng ghế. Hàng ghế đầu tiên (gần sân khấu nhất) có $11$ ghế, mỗi hàng ghế phía sau nhiều hơn hàng liền trước đúng $3$ ghế. Tính tổng số ghế của khu vực khán giả.}{1261}

\cauTLN[id=T11-C2-B6-D10-VD-090,muc=VD,loigiai={Tìm $n$ lớn nhất với $\dfrac{n[2\cdot4+(n-1)\cdot2]}{2}\le126$, được $n=9$.}]{Người ta xếp các thùng sơn thành một khối hình thang trong kho: hàng trên cùng có $4$ thùng, mỗi hàng bên dưới nhiều hơn hàng liền trên $2$ thùng. Có $126$ thùng sơn. Hỏi xếp được nhiều nhất bao nhiêu hàng đầy đủ theo quy luật trên?}{9}

\cauTN[id=T11-C2-B6-D10-VDC-091,muc=VDC,loigiai={Gọi ba số là $a-d$, $a$, $a+d$ ($d>0$): $3a=21\Rightarrow a=7$; $(7-d)\cdot7\cdot(7+d)=168\Rightarrow d=5$.}]{Ba số lập thành một cấp số cộng tăng có tổng bằng $21$ và tích bằng $168$. Số lớn nhất trong ba số đó là}
  {$12$.}
  {$14$.}{$11$.}{$2$.}

\cauTN[id=T11-C2-B6-D10-VDC-092,muc=VDC,loigiai={Lương tháng các năm lập thành cấp số cộng $u_1=7$, $d=\dfrac{1}{2}$; tổng $=12S_{5}=480$ triệu đồng.}]{Công ty TNHH Đại Việt tuyển kĩ sư với mức lương năm đầu tiên là $7$ triệu đồng/tháng. Theo hợp đồng, sau mỗi năm làm việc, mức lương mỗi tháng được tăng thêm $\dfrac{1}{2}$ triệu đồng so với năm liền trước. Tổng số tiền lương (triệu đồng) kĩ sư nhận được sau $5$ năm làm việc là}
  {$480$.}
  {$486$.}{$40$.}{$479$.}

\cauTN[id=T11-C2-B6-D10-VDC-093,muc=VDC,loigiai={$S_{22}=\dfrac{22(2\cdot23+21\cdot3)}{2}=1199$ ghế.}]{Khu vực ghế khán giả của rạp chiếu phim của một trung tâm thương mại được bố trí theo hình rẻ quạt gồm $22$ hàng ghế. Hàng ghế đầu tiên (gần sân khấu nhất) có $23$ ghế, mỗi hàng ghế phía sau nhiều hơn hàng liền trước đúng $3$ ghế. Tổng số ghế của khu vực khán giả là}
  {$1199$.}
  {$1958$.}{$1268$.}{$1176$.}

\cauTN[id=T11-C2-B6-D10-VDC-094,muc=VDC,loigiai={Gọi ba số là $a-d$, $a$, $a+d$ ($d>0$): $3a=9\Rightarrow a=3$; $(3-d)\cdot3\cdot(3+d)=-21\Rightarrow d=4$.}]{Ba số lập thành một cấp số cộng tăng có tổng bằng $9$ và tích bằng $-21$. Số lớn nhất trong ba số đó là}
  {$7$.}
  {$8$.}{$3$.}{$-1$.}

\cauTN[id=T11-C2-B6-D10-VDC-095,muc=VDC,loigiai={$S_{20}=\dfrac{20(2\cdot17+19\cdot2)}{2}=720$ ghế.}]{Khu vực ghế khán giả của nhà văn hoá thiếu nhi của tỉnh được bố trí theo hình rẻ quạt gồm $20$ hàng ghế. Hàng ghế đầu tiên (gần sân khấu nhất) có $17$ ghế, mỗi hàng ghế phía sau nhiều hơn hàng liền trước đúng $2$ ghế. Tổng số ghế của khu vực khán giả là}
  {$720$.}
  {$762$.}{$719$.}{$722$.}

\cauTN[id=T11-C2-B6-D10-VDC-096,muc=VDC,loigiai={Gọi ba số là $a-d$, $a$, $a+d$ ($d>0$): $3a=21\Rightarrow a=7$; $(7-d)\cdot7\cdot(7+d)=231\Rightarrow d=4$.}]{Ba số lập thành một cấp số cộng tăng có tổng bằng $21$ và tích bằng $231$. Số lớn nhất trong ba số đó là}
  {$11$.}
  {$3$.}{$12$.}{$15$.}

\cauDS[id=T11-C2-B6-D10-VDC-097,muc=VDC,loigiai={Cấp số cộng $u_1=12$, $d=3$: $u_{15}=54$, $u_{16}=57$, $S_{16}=552$. a) Đúng. b) Đúng. c) Đúng. d) Sai.}]{Khu vực ghế khán giả của Nhà hát Lớn Hà Nội được bố trí theo hình rẻ quạt gồm $16$ hàng ghế. Hàng ghế đầu tiên (gần sân khấu nhất) có $12$ ghế, mỗi hàng ghế phía sau nhiều hơn hàng liền trước đúng $3$ ghế.}{%
  \y{Hàng thứ hai có $15$ ghế.}{d}
  \y{Hàng thứ $15$ có $54$ ghế.}{d}
  \y{Hàng cuối cùng có $57$ ghế.}{d}
  \y{Khu vực khán giả có tất cả $564$ ghế.}{s}}

\cauDS[id=T11-C2-B6-D10-VDC-098,muc=VDC,loigiai={Cấp số cộng $u_1=21$, $d=2$: $u_{13}=45$, $u_{26}=71$, $S_{26}=1196$. a) Đúng. b) Đúng. c) Đúng. d) Sai.}]{Khu vực ghế khán giả của hội trường của trường THPT Nguyễn Trãi được bố trí theo hình rẻ quạt gồm $26$ hàng ghế. Hàng ghế đầu tiên (gần sân khấu nhất) có $21$ ghế, mỗi hàng ghế phía sau nhiều hơn hàng liền trước đúng $2$ ghế.}{%
  \y{Hàng thứ hai có $23$ ghế.}{d}
  \y{Hàng thứ $13$ có $45$ ghế.}{d}
  \y{Hàng cuối cùng có $71$ ghế.}{d}
  \y{Khu vực khán giả có tất cả $1217$ ghế.}{s}}

\cauTLN[id=T11-C2-B6-D10-VDC-099,muc=VDC,loigiai={Giải $\dfrac{n[2\cdot(-6)+(n-1)\cdot4]}{2}=1242$, $n\in\mathbb{N}^*$ được $n=27$.}]{Cho cấp số cộng $(u_n)$ có $u_1=-6$, $d=4$. Biết tổng $n$ số hạng đầu bằng $1242$. Tìm $n$.}{27}

\cauTLN[id=T11-C2-B6-D10-VDC-100,muc=VDC,loigiai={Tìm $n$ lớn nhất với $\dfrac{n[2\cdot7+(n-1)\cdot1]}{2}\le287$, được $n=18$.}]{Người ta xếp các thùng sơn thành một khối hình thang trong kho: hàng trên cùng có $7$ thùng, mỗi hàng bên dưới nhiều hơn hàng liền trên $1$ thùng. Có $287$ thùng sơn. Hỏi xếp được nhiều nhất bao nhiêu hàng đầy đủ theo quy luật trên?}{18}
